\section*{ВВЕДЕНИЕ}
\addcontentsline{toc}{section}{ВВЕДЕНИЕ}

В современном мире йога как практика физического и духовного развития становится всё более популярной. Миллионы людей занимаются йогой для поддержания здоровья, снятия стресса и улучшения качества жизни. Однако многие йога-студии сталкиваются с проблемами автоматизации процесса записи клиентов на занятия, управления расписанием и отслеживания прогресса.

Традиционные методы управления (запись по телефону, журналы, таблицы Excel) являются неэффективными и отнимают много времени как у администраторов, так и у клиентов. Клиенты хотят иметь возможность записаться на занятие в любое время, получить напоминание и отслеживать свой прогресс. Администраторам необходим инструмент для управления расписанием, инструкторами и просмотра всех записей.

\textbf{Актуальность} данной работы обусловлена необходимостью создания удобного, бесплатного и автономного инструмента для записи на занятия йогой, который мог бы использоваться как отдельными студиями, так и конечными пользователями.

\textbf{Объектом исследования} являются процессы автоматизации записи на занятия в йога-студиях.

\textbf{Предметом исследования} — методы и средства разработки мобильных приложений для организации занятий йогой.

\textbf{Целью данной работы} является разработка программно-информационной системы для организации и проведения занятий йогой в виде мобильного приложения для платформы Android.

Для достижения поставленной цели необходимо решить следующие задачи:
\begin{enumerate}
	\item провести анализ предметной области и существующих аналогов;
	\item сформулировать функциональные и нефункциональные требования к разрабатываемой системе;
	\item обосновать выбор технологий разработки;
	\item разработать архитектуру приложения;
	\item спроектировать пользовательский интерфейс;
	\item реализовать основные модули приложения;
	\item выполнить тестирование разработанного приложения.
\end{enumerate}

\textbf{Методы исследования:} анализ научной и технической литературы, сравнительный анализ существующих решений, моделирование с использованием UML, объектно-ориентированное проектирование, тестирование программного обеспечения.

\textbf{Научная новизна} заключается в разработке автономного мобильного приложения для записи на занятия йогой, не требующего подключения к интернету и использования серверной инфраструктуры.

\textbf{Практическая значимость} работы состоит в том, что разработанное приложение может быть использовано йога-студиями и фитнес-центрами для автоматизации записи клиентов, а также отдельными пользователями для отслеживания личного прогресса.

Результатом работы является мобильное приложение «Йога-студия» для платформы Android, позволяющее пользователям просматривать расписание, записываться на занятия, отслеживать прогресс, оставлять отзывы и получать уведомления.